\section{Differentiation and integration of power series}
\label{sec:power-calculus-new}
\textbf{The permission comes before the calculation.}
Inside the radius of convergence of a power series, we may differentiate
term by term, and we may integrate term by term from zero to any interior
point. Uniform convergence on each smaller closed interval is the reason
integration works. Endpoints require a separate argument: convergence inside
an interval does not by itself settle either endpoint. We use Abel's theorem
only after checking that the numerical endpoint series converges.

\exercise{26.1}{Term-by-term integration at a positive input}
\goal If $f(t)=\sum_{k=0}^\infty a_kt^k$ has radius $R>0$, prove
$\int_0^x f(t)\,dt=\sum_{k=0}^\infty a_kx^{k+1}/(k+1)$ when $0<x<R$.
\step{Work on an interval lying strictly inside the radius.}
Fix such an $x$, and put $s_n(t)=\sum_{k=0}^n a_kt^k$.
The power series converges uniformly on $[-x,x]$ by Theorem~26.1, and
therefore also on $[0,x]$. Each $s_n$ is a continuous polynomial.
\step{Pass the limit through the integral using the established theorem.}
Theorem~25.2 yields
\[
 \int_0^x f(t)\,dt=\lim_{n\to\infty}\int_0^x s_n(t)\,dt.
\]
For a finite sum, linearity and the power rule give
\[
 \int_0^x s_n(t)\,dt
 =\sum_{k=0}^n a_k\int_0^x t^k\,dt
 =\sum_{k=0}^n \frac{a_k}{k+1}x^{k+1}.
\]
Taking the limit of these finite sums gives precisely the claimed infinite series.
\answer{The term-by-term integration formula holds for every positive $x<R$.
The strict inequality $x<R$ is what secured the uniform convergence used in the proof.}

\exercise{26.2}{A differentiated geometric series and three numerical sums}
\goal Derive $\sum_{n=1}^\infty nx^n$ and use it to evaluate the requested sums.
\step{(a) Differentiate first, then restore the missing power of $x$.}
For $|x|<1$, the geometric series is
$\sum_{n=0}^\infty x^n=1/(1-x)$.
Its radius is $1$, so Theorem~26.5 allows differentiation at every such $x$:
\[
 \sum_{n=1}^\infty nx^{n-1}=\frac1{(1-x)^2}.
\]
Multiplying by $x$ gives
\[
 \boxed{\sum_{n=1}^\infty nx^n=\frac{x}{(1-x)^2}}
 \qquad(|x|<1).
\]
\step{(b) Substitute $x=1/2$.}
The value lies strictly inside the radius, so the formula gives
\[
 \sum_{n=1}^\infty\frac n{2^n}
 =\frac{1/2}{(1-1/2)^2}=\frac{1/2}{1/4}=2.
\]
\step{(c) Substitute the two signed inputs separately.}
At $x=1/3$,
\[
 \sum_{n=1}^\infty\frac n{3^n}
 =\frac{1/3}{(2/3)^2}=\frac34.
\]
At $x=-1/3$,
\[
 \sum_{n=1}^\infty\frac{(-1)^n n}{3^n}
 =\frac{-1/3}{(1+1/3)^2}
 =\frac{-1/3}{16/9}=-\frac3{16}.
\]
\answer{The three numerical sums are $2$, $3/4$, and $-3/16$.}

\exercise{26.3}{One more derivative gives squared indices}
\goal Find $\sum_{n=1}^\infty n^2x^n$ and evaluate it at $1/2$ and $1/3$.
\step{(a) Differentiate the identity from Exercise~26.2.}
For $|x|<1$, term-by-term differentiation is legitimate, because differentiation
preserves the radius of a power series. Hence
\[
 \sum_{n=1}^\infty n^2x^{n-1}
 =\frac{d}{dx}\left[\frac{x}{(1-x)^2}\right]
 =\frac1{(1-x)^2}+\frac{2x}{(1-x)^3}
 =\frac{1+x}{(1-x)^3}.
\]
Multiplication by $x$ gives
\[
 \boxed{\sum_{n=1}^\infty n^2x^n=\frac{x(1+x)}{(1-x)^3}}
 \qquad(|x|<1).
\]
\step{(b) Evaluate the two substitutions.}
Both inputs are strictly inside the radius. At $x=1/2$ the numerator is
$(1/2)(3/2)=3/4$ and the denominator is $(1/2)^3=1/8$, so
\[
 \sum_{n=1}^\infty\frac{n^2}{2^n}=\frac{3/4}{1/8}=6.
\]
At $x=1/3$ the numerator is $4/9$ and the denominator is $8/27$, so
\[
 \sum_{n=1}^\infty\frac{n^2}{3^n}=\frac{4/9}{8/27}=\frac32.
\]
\answer{The general formula is $x(1+x)/(1-x)^3$, and the requested values
are $6$ and $3/2$.}

\exercise{26.4}{Integrating the Gaussian power series}
\goal Express $e^{-x^2}$ and $F(x)=\int_0^x e^{-t^2}\,dt$ as power series.
\step{(a) Substitute into the exponential series.}
The exponential identity, with its constant term included, is
\[
 e^u=\sum_{n=0}^\infty\frac{u^n}{n!}\qquad(u\in\R).
\]
Taking $u=-x^2$ gives
\[
 e^{-x^2}=\sum_{n=0}^\infty\frac{(-1)^n}{n!}x^{2n}
 =1-x^2+\frac{x^4}{2!}-\frac{x^6}{3!}+\cdots.
\]
This substitution is valid for every real $x$ because the original exponential
series converges for every real input.
\step{(b) Verify the radius before integrating.}
For fixed $x$, the ratio of successive absolute terms is $x^2/(n+1)\to0$,
so this power series has infinite radius. Theorem~26.4 therefore applies
between zero and every real $x$, whether positive or negative. It gives
\[
 \boxed{F(x)=\sum_{n=0}^\infty
       \frac{(-1)^n x^{2n+1}}{(2n+1)n!}}
 =x-\frac{x^3}{3}+\frac{x^5}{10}-\frac{x^7}{42}+\cdots.
\]
The lower limit contributes zero for each monomial, so no integration constant
has been omitted.
\answer{The two displayed expansions hold for all $x\in\R$.}

\exercise{26.5}{A function equal to its own derivative, without naming it}
\goal For $f(x)=\sum_{n=0}^\infty x^n/n!$, prove $f'=f$ directly from the series.
\step{Check that term-by-term differentiation is allowed everywhere.}
For each fixed $x$, the ratio of successive absolute terms is
$|x|/(n+1)\to0$. Thus the radius is infinite, and Theorem~26.5 allows
term-by-term differentiation at every real input.
\step{Differentiate and carefully change the index.}
The constant term differentiates to zero, and for $n\ge1$,
$n/n!=1/(n-1)!$. Therefore
\[
 f'(x)=\sum_{n=1}^\infty\frac{nx^{n-1}}{n!}
 =\sum_{n=1}^\infty\frac{x^{n-1}}{(n-1)!}
 =\sum_{j=0}^\infty\frac{x^j}{j!}=f(x).
\]
In the last step we set $j=n-1$, so the first new index is zero.
\answer{$f'=f$ on $\R$. No identification of $f$ with the exponential
function was used.}

\exercise{26.6}{A trigonometric identity derived without trigonometry}
\goal For
\[
 s(x)=\sum_{n=0}^\infty\frac{(-1)^n x^{2n+1}}{(2n+1)!},\qquad
 c(x)=\sum_{n=0}^\infty\frac{(-1)^n x^{2n}}{(2n)!},
\]
prove the requested derivative relations and $s^2+c^2=1$.
\step{(a) Differentiate each series with its radius justified.}
For fixed $x$, ratios of successive absolute terms tend to zero in both series,
since their denominators acquire two factorial factors while their numerators
acquire only a factor $x^2$. Thus both radii are infinite.
Theorem~26.5 gives
\[
 s'(x)=\sum_{n=0}^\infty\frac{(-1)^n x^{2n}}{(2n)!}=c(x).
\]
For $c$, the $n=0$ term differentiates to zero, so
\begin{align*}
 c'(x)&=\sum_{n=1}^\infty\frac{(-1)^n x^{2n-1}}{(2n-1)!}\\
 &=-\sum_{j=0}^\infty\frac{(-1)^j x^{2j+1}}{(2j+1)!}=-s(x).
\end{align*}
Here $j=n-1$, and $(-1)^{j+1}=-(-1)^j$ explains the minus sign.
\step{(b) Differentiate the sum of squares.}
The product rule gives
\[
 (s^2+c^2)'=2ss'+2cc'=2sc+2c(-s)=0.
\]
\step{(c) Determine the constant from one input.}
A differentiable function with derivative zero on an interval is constant
there, by the Mean Value theorem. Apply this on the interval joining zero to
any real $x$. Thus $s(x)^2+c(x)^2=s(0)^2+c(0)^2$.
The defining series give $s(0)=0$ and $c(0)=1$, so that constant is $1$.
\answer{$s'=c$, $c'=-s$, and $s(x)^2+c(x)^2=1$ for every real $x$,
proved using the series alone.}

\exercise{26.7}{Why absolute value has no power-series representation at zero}
\goal Decide whether $|x|=\sum_{n=0}^\infty a_nx^n$ can hold for all real $x$.
\step{A power series with positive radius is differentiable at its center.}
If such a representation held for every real input, its radius would be infinite,
in particular positive. Theorem~26.5 would make its sum differentiable at zero.
\step{Calculate the two difference quotients of absolute value.}
Since $|0|=0$,
\[
 \frac{|h|-|0|}{h}=\begin{cases}1&h>0,\\-1&h<0.\end{cases}
\]
The right-hand and left-hand limits differ, so $|x|$ is not differentiable at zero.
This contradicts the consequence of a power-series representation.
\answer{No such power series exists, even on any two-sided neighborhood of
zero. On the separate half-lines the formulas are $|x|=x$ and $|x|=-x$, but
these do not combine into one power series about zero. Polynomial
approximation, discussed next, is a different and weaker requirement.}

\exercise{26.8}{The arctangent series and its endpoints}
\goal Derive the arctangent series inside $(-1,1)$ and settle both endpoints.
\step{(a) Substitute into the geometric series.}
For $|x|<1$, put $y=-x^2$, so $|y|<1$. Then
\[
 \sum_{n=0}^\infty(-1)^n x^{2n}
 =\sum_{n=0}^\infty y^n=\frac1{1-y}=\frac1{1+x^2}.
\]
\step{(b) Integrate on a strictly interior interval.}
The power series has radius $1$, so Theorem~26.4 gives, for $|x|<1$,
\[
 \int_0^x\frac{dt}{1+t^2}
 =\sum_{n=0}^\infty\frac{(-1)^n x^{2n+1}}{2n+1}.
\]
The integral on the left is $\arctan x-\arctan0=\arctan x$. Therefore
\[
 \arctan x=x-\frac{x^3}{3}+\frac{x^5}{5}-\cdots
 \qquad(-1<x<1).
\]
\step{(c) Establish convergence at $1$ before using Abel's theorem.}
At $x=1$ the numerical series is
$\sum_{n=0}^\infty(-1)^n/(2n+1)$. Its positive term sizes decrease to zero,
so the alternating-series test proves convergence.
Abel's theorem then identifies its sum with the limit of the interior sum
as $x\uparrow1$. Since arctangent is continuous,
\[
 \sum_{n=0}^\infty\frac{(-1)^n}{2n+1}
 =\lim_{x\uparrow1}\arctan x=\frac\pi4.
\]
Multiplication by $4$ gives the familiar formula
\[
 \boxed{\pi=4\left(1-\frac13+\frac15-\frac17+\cdots\right)}.
\]
\step{(d) Evaluate at the negative endpoint with its sign intact.}
Every power in the series is odd. Substitution $x=-1$ therefore makes the
series the negative of the convergent series in part (c). Its sum is
$-\pi/4=\arctan(-1)$. It converges conditionally at both endpoints: the
series of absolute values $\sum1/(2n+1)$ diverges by comparison with the
harmonic series, since $1/(2n+1)\ge1/(2n+2)$.
\answer{The arctangent expansion holds on $[-1,1]$, with endpoint values
$\pi/4$ and $-\pi/4$ justified separately, rather than by interior convergence alone.}
